\section{Análisis de casos: tareas de Web, código e investigación}

\subsection{Caso uno: interacción de cadena larga del Web Agent}
En escenarios de Web, el Agent suele necesitar navegar entre páginas, analizar información semiestructurada y completar decisiones de varios pasos. El punto de vista del curso es que el fracaso de este tipo de tareas normalmente no se debe a que «el modelo no entiende la página», sino a que «la sincronización de estado y la actualización del plan no son oportunas». Los problemas típicos incluyen: cambios en el DOM de la página que invalidan la localización, interrupciones del estado de inicio de sesión, palabras objetivo sinónimas que provocan clics erróneos, y valores devueltos por las herramientas que no coinciden con el estado visual.

\begin{importantbox}{La dificultad central del Web Agent es «un control de circuito cerrado sólido en un entorno de estructura débil»}
El Agent debe completar de manera continua el circuito cerrado Observe-Reason-Act bajo un estado de interfaz inestable. Si carece de una máquina de estados explícita y de verificadores a nivel de paso, los errores se propagarán entre turnos y terminarán manifestándose como fallos aleatorios.
\end{importantbox}

Una estrategia viable es dividir las operaciones de página en tres categorías: «acciones de navegación», «acciones de extracción» y «acciones de confirmación», y configurar una lógica de verificación distinta para cada tipo. Por ejemplo, las acciones de navegación se centran en los cambios de URL o título de página, las acciones de extracción se centran en la integridad de los campos, y las acciones de confirmación se centran en si se cumplen las restricciones de negocio. Esto permite reducir la probabilidad de que «un paso erróneo arruine todo el proceso».

\begin{knowledgebox}{Recomendaciones de evaluación para el Web Agent}
\begin{itemize}
    \item Registrar la tasa de éxito a nivel de acción, no solo la tasa de éxito del resultado final;
    \item Contabilizar el número de reintentos y de retrocesos para evaluar la estabilidad del plan;
    \item Agrupar por complejidad de página, para evitar que el promedio oculte los casos difíciles.
\end{itemize}
\end{knowledgebox}

\subsection{Caso dos: tareas a nivel de repositorio del Code Agent}
Las tareas de código parecen estructuradas, pero los objetivos a nivel de repositorio suelen involucrar dependencias entre archivos, convenciones implícitas y restricciones históricas. El curso enfatiza que la vía de mejora del Code Agent no es solo «una generación de código más fuerte», sino «una comprensión más fuerte del estado del repositorio + un ciclo de retroalimentación de verificación más fuerte».

\begin{warningbox}{Hacer solo parches de una ronda amplifica la deuda técnica}
Si el Agent solo persigue que las pruebas actuales pasen, e ignora la consistencia de la arquitectura, las reglas de nomenclatura y los contratos de interfaz, el éxito a corto plazo se convertirá en un costo de mantenimiento a largo plazo. El Code Agent necesita usar lint, type check, unit test e integration test en conjunto como verificador.
\end{warningbox}

A nivel de implementación, la tarea puede dividirse en cuatro etapas: «comprensión», «modificación», «verificación» y «revisión posterior», y la salida de cada etapa se escribe en la episodic memory. El valor de hacerlo así es que, cuando la tarea falla, el sistema puede ubicar si se trata de una desviación en la comprensión de los requisitos, un error de implementación o una cobertura de verificación insuficiente, y así activar una replanificación específica.

\begin{knowledgebox}{Recomendaciones de memory para el Code Agent}
\begin{itemize}
    \item semantic memory almacena las convenciones del proyecto, las interfaces públicas y las decisiones históricas;
    \item episodic memory almacena los intentos clave y las razones de fracaso de esta tarea;
    \item procedural memory almacena scripts reutilizables (como flujos de trabajo de pruebas de regresión).
\end{itemize}
\end{knowledgebox}

\subsection{Caso tres: el circuito cerrado de evidencia del Research Agent}
Las tareas de investigación a menudo se juzgan erróneamente como «puros problemas de reasoning». Dentro del marco de Yu Su, el Research Agent también necesita un memory y un planning sólidos: el primero se usa para la consolidación de evidencia y la gestión de conflictos, y el segundo para el cambio de ruta de investigación y la asignación del presupuesto experimental.

\begin{importantbox}{Lo clave del Research Agent no es «escribir como un artículo académico», sino «que la evidencia sea rastreable»}
El sistema debe vincular cada conclusión con su fuente, su nivel de confianza y su marca de tiempo. De lo contrario, cuando el corpus de recuperación se actualice o cambien los valores devueltos por las herramientas, las conclusiones antiguas seguirán contaminando el razonamiento posterior dentro de la capa de memoria.
\end{importantbox}

Una práctica útil es adoptar una estructura ternaria de «conclusión-evidencia-contraevidencia» al escribir en la memory: cada vez que el Agent presenta una conclusión, debe adjuntar evidencia de respaldo y posibles contraejemplos. Esto puede mejorar de manera significativa la eficacia de la etapa de reflexión (reflection) y reducir las «alucinaciones de tipo explicativo».

\begin{table}[H]
\centering
\begin{tabular}{p{2.7cm}p{3.9cm}p{4.4cm}}
\toprule
\textbf{Tipo de tarea} & \textbf{Capacidad prioritaria} & \textbf{Riesgo principal} \\
\midrule
Web Agent & Sincronización de estado + replanificación dinámica & La deriva de página invalida las acciones \\
Code Agent & Circuito cerrado de verificación + memoria del repositorio & Los parches localmente óptimos erosionan la calidad global \\
Research Agent & Gestión de evidencia + mecanismo de reflexión & Contaminación de la memoria por conclusiones sin fuente \\
\bottomrule
\end{tabular}
\caption{Diferencias en la prioridad de capacidades entre los tres tipos de tareas representativas}
\end{table}

\subsection{Resumen de la sección}
El punto en común de los tres tipos de tareas es que todas requieren un circuito cerrado del sistema; la diferencia radica en que el enfoque de diseño del verificador y de la estrategia de memoria es distinto. El desarrollo del Agent debe configurar las capacidades según los atributos de la tarea, en lugar de perseguir una plantilla única y genérica.
